\section{Verification observables and actual refinement evidence}
\label{sec:verification}
\subsection{What the fixture observables measure}
The particle spacing at initialization is $s_p=h/2$. The reported surface
proxies are
\begin{equation}
 H_h=\max_p z_p+s_p/2,\qquad
 B_{x,h}=\max_p x_p-\min_p x_p+s_p,\qquad S_h=H_0-H_h.
 \label{eq:surface-proxies}
\end{equation}
These estimate height, horizontal extent and slump using particle extrema plus
an \emph{initial} spacing correction. They are not reconstructed deformed
particle surfaces or a continuum free-surface mesh. They can therefore have
quadrature and surface-estimation error in addition to solver error.

The hydrostatic error metrics compare every particle, including those next to
the bottom boundary, to the phase-appropriate reference in
Section~\ref{sec:constitutive}:
\begin{equation}
 e_{\rm rms}=\left[\frac1{N_p}\sum_p
 \left(\frac{q_p-q_{\rm ref}(z_p)}{\rho_0gH_0}\right)^2\right]^{1/2},
 \qquad e_{\max}=\max_p\frac{|q_p-q_{\rm ref}(z_p)|}{\rho_0gH_0}.
\end{equation}
For liquid $q=p$ is pressure; for paste it is negative vertical Kirchhoff stress.
Removing bottom-row particles or smoothing per-particle liquid pressure would
change the acceptance question, not solve the observed accuracy problem.
The archived dataset below contains slump results only; it establishes no
hydrostatic-pressure convergence claim.

The code also reports a cylinder-slump sanity estimate \cite{pashias},
\begin{equation}
 Y=\frac{\tau_y}{\rho_0gH_0},\qquad
 \frac{S_{\rm cyl}}{H_0}=1-2Y[1-\ln(2Y)],\quad 0<Y<\tfrac12.
\end{equation}
It uses zero slump for $Y\geq1/2$ and $H_0$ for $Y\leq0$.
For $\tau_y=\SI{30}{Pa}$, $\rho_0=\SI{1000}{kg.m^{-3}}$,
$H_0=\SI{20}{mm}$ and $g=\SI{9.81}{m.s^{-2}}$, $Y\simeq0.152905$ and
$S_{\rm cyl}\simeq\SI{6.63737}{mm}$. The simulated block is a \emph{cube},
with finite elasticity, wall friction and time-dependent rheology. A mismatch
with this cylinder estimate is not by itself a validated error measure for
that different problem.

\subsection{Separate spatial and temporal experiments}
For observable $Q$ at three or more successively finer resolutions, the
implemented criterion computes
\begin{equation}
 \epsilon_k=\frac{|Q_k-Q_{k-1}|}{\max(|Q_k|,Q_{\rm scale})},\qquad
 \epsilon_{\rm final}<0.05. \label{eq:refinement}
\end{equation}
A finest value below $Q_{\rm scale}$ is \emph{uninformative}, not a pass. The
threshold is strict. Missing/nonfinite values, incomplete runs, changed fidelity,
wrong final times, adaptive limiting, rejected steps and budget overruns prevent
passing claims. The spatial axis holds the requested timestep fixed; the temporal
axis holds $h$ fixed and requires increasing actual step counts.

Three values can expose a failure to resolve an observable, but a small final
change does not prove asymptotic convergence. For equal refinement ratio $r$,
\begin{equation}
 p_{\rm obs}=\frac{\ln\left(|Q_0-Q_1|/|Q_1-Q_2|\right)}{\ln r}
\end{equation}
is merely a trend diagnostic without a verified leading-error expansion.
Neither Richardson extrapolation nor a formal discretization uncertainty bound
is claimed by this study.

The frozen run uses \path{examples/research_slump.yaml}, modified only by the
study axes in \path{examples/verification_slump.yaml}: a \SI{20}{mm} cube in a
\SI{60}{mm} domain, $\rho_0=\SI{1000}{kg.m^{-3}}$, $E=\SI{1000}{Pa}$,
$\nu_P=0.2$, $\tau_y=\SI{30}{Pa}$, $K_s=\SI{10}{Pa.s^{0.6}}$,
$n=0.6$, seed 7, final time \SI{0.5}{s}. All material values are synthetic.
The undeformed block mass is \SI{\BlockMassGrams}{g}, with continuum initial
potential energy $mgH_0/2=\SI{\BlockPotentialMillijoules}{mJ}$.
Non-hydrostatic fixtures add bounded position jitter of amplitude $0.1s_p$;
therefore the discrete initial potential energy differs slightly from this
integral. Spatial refinement also refines particle quadrature and jitter
amplitude; it is not a grid-only experiment with fixed material points.

\begin{table}[htbp]
\centering\small
\begin{tabular}{lrrrrrr}
\toprule
Axis & $h$ (mm) & $\Delta t$ (ms) & Points & Steps & Slump (mm) & Spread (mm) \\
\midrule
Spatial & 10 & 0.05 & 64 & 10000 & 3.9897 & 22.6655 \\
Spatial & 5 & 0.05 & 512 & 10000 & 4.3623 & 23.5512 \\
Spatial & 2.5 & 0.05 & 4096 & 10000 & 4.5165 & 25.9463 \\
Temporal & 5 & 0.4 & 512 & 1250 & 4.7376 & 25.0541 \\
Temporal & 5 & 0.2 & 512 & 2500 & 4.6077 & 24.6164 \\
Temporal & 5 & 0.1 & 512 & 5000 & 4.4875 & 24.1053 \\
\bottomrule
\end{tabular}

\caption{Actual compiled-kernel results at \SI{0.5}{s}, source revision
\nolinkurl{9a66a55}. Every case completed with zero limited and rejected steps.
Times and lengths here are requested caps/grid scales and measured final
surface proxies, respectively.}
\label{tab:cases}
\end{table}

\begin{figure}[p]
\centering
\begin{tikzpicture}
\begin{groupplot}[group style={group size=3 by 2,horizontal sep=1.5cm,vertical sep=2cm},
 width=.295\textwidth,height=.25\textheight,
 xmode=log,log basis x=2,x dir=reverse,grid=major,
 tick label style={font=\footnotesize},label style={font=\small},
 title style={font=\small},scaled ticks=false]
\nextgroupplot[title={Slump},xlabel={$h$ (mm)},ylabel={mm},xtick={10,5,2.5},
 xticklabels={10,5,2.5}]
\addplot+[blue!75!black,mark=*,thick] table[x=h_mm,y=slump_mm,col sep=comma]{generated/spatial.csv};
\nextgroupplot[title={Horizontal spread},xlabel={$h$ (mm)},ylabel={mm},xtick={10,5,2.5},
 xticklabels={10,5,2.5}]
\addplot+[blue!75!black,mark=*,thick] table[x=h_mm,y=spread_mm,col sep=comma]{generated/spatial.csv};
\nextgroupplot[title={Plastic-loss estimate},xlabel={$h$ (mm)},ylabel={$\mu$J},xtick={10,5,2.5},
 xticklabels={10,5,2.5}]
\addplot+[blue!75!black,mark=*,thick] table[x=h_mm,y=plastic_uj,col sep=comma]{generated/spatial.csv};
\nextgroupplot[xlabel={$\Delta t$ (ms)},ylabel={mm},xtick={.4,.2,.1},
 xticklabels={0.4,0.2,0.1}]
\addplot+[orange!85!black,mark=square*,thick] table[x=dt_ms,y=slump_mm,col sep=comma]{generated/temporal.csv};
\nextgroupplot[xlabel={$\Delta t$ (ms)},ylabel={mm},xtick={.4,.2,.1},
 xticklabels={0.4,0.2,0.1}]
\addplot+[orange!85!black,mark=square*,thick] table[x=dt_ms,y=spread_mm,col sep=comma]{generated/temporal.csv};
\nextgroupplot[xlabel={$\Delta t$ (ms)},ylabel={$\mu$J},xtick={.4,.2,.1},
 xticklabels={0.4,0.2,0.1}]
\addplot+[orange!85!black,mark=square*,thick] table[x=dt_ms,y=plastic_uj,col sep=comma]{generated/temporal.csv};
\end{groupplot}
\end{tikzpicture}
\caption{Frozen simulation evidence, \emph{not} illustrative curves. Top row:
spatial refinement at fixed $\Delta t=\SI{0.05}{ms}$. Bottom row: temporal
refinement at fixed $h=\SI{5}{mm}$. Each axis becomes finer to the right;
lines connect only the three observed values and are not fitted convergence
curves. Vertical axes are independent between panels. The study remains
unresolved because spread and plastic loss fail the specified final-two tests.}
\label{fig:refinement}
\end{figure}

\begin{table}[htbp]
\centering
\begin{tabular}{lrrr}
\toprule
Final-two change & Slump & Spread & Plastic loss \\
\midrule
Spatial & 3.41\% & 9.23\% & 31.39\% \\
Temporal & 2.68\% & 2.12\% & 10.25\% \\
\bottomrule
\end{tabular}

\caption{Final-two relative changes computed with \eqref{eq:refinement}.
Required: each listed value strictly below 5\%. Near-zero scales are
\SI{0.5}{mm} for lengths and \SI{1}{\micro J} for plastic loss.}
\label{tab:verdicts}
\end{table}

At $h=\SI{5}{mm}$ and $\Delta t=\SI{0.1}{ms}$, the recorded values include
\begin{align}
 E_N&=-\SI{109.379}{\micro J}, &D_F&=\SI{35.486}{\micro J},\\
 D_P&=\SI{173.886}{\micro J}, &r_E&=\SI{84.476}{\micro J}.
\end{align}
The nonzero diagnostic is retained. It is neither a pass of energy closure nor
an external-wall-work measurement. The six-case study took approximately
\SI{11.86}{s} on the producing arm64, 10-core, 16 GiB machine, measured through
report evaluation (not the final report write). This is a tiny-fixture timing,
not evidence for full-box or multi-day feasibility.

\paragraph{Implementation mapping and evidence identity.}
\path{core/src/mpm/fixtures.rs}: \nolinkurl{FixtureSpec::shape_observables},
\nolinkurl{slump_observables}, \nolinkurl{hydrostatic_observables}.
\path{python/litter_physics/verification.py}: \nolinkurl{plan_cases},
\nolinkurl{evaluate_observable}, \nolinkurl{evaluate_axis}, \nolinkurl{run_study}.
The verbatim archived report is \path{docs/formal/data/refinement-report.json};
its source commit, clean-tree flag, native/Python hashes, dependency versions,
hardware, all case values and unresolved verdicts are retained. The manuscript
build verifies the archive's SHA-256 before extracting any plot or table data.

\section{Calibration, reduced responses, and uncertainty}
\label{sec:calibration}
\subsection{What can currently be fitted}
The Python calibration module contains bounded empirical fits including
\begin{equation}
 U(t)=C[1-e^{-k_ut}],\qquad B(t)=1-e^{-k_bt},\qquad
 \tau(\dot\gamma)=\tau_y+K_s\dot\gamma^n.
\end{equation}
Tilt observations supply a friction prior based on $\tan\theta$; they do not
uniquely identify all multisphere sliding, rolling and shape parameters.
Visit observations fit stage statistics and can be held out for personalization
scoring. Synthetic observations remain explicitly synthetic; fit recovery from
synthetic data cannot identify a real cat or material.

There is an important model discrepancy: the fitted breakup fraction $B(t)$
is exponential, whereas the household runtime uses the accumulated,
whole-pellet threshold in \eqref{eq:damage}. Transferring $k_b$ between them does
not automatically calibrate the distribution of actual breakup times. A
simulator-based likelihood, an explicit heterogeneous threshold population, or
measured justification of the reduction is still needed. Likewise, fitting a
flow curve alone does not identify bulk elasticity, adhesion or porosity.

A prospective joint inverse problem can be written
\begin{equation}
 \widehat{\boldsymbol\theta}=\arg\min_{\boldsymbol\theta\in\Theta}
 \sum_j\frac{[\mathcal H_j(\mathcal S(\boldsymbol\theta))-y_j]^2}
                 {s_j^2}+\mathcal R(\boldsymbol\theta), \label{eq:inverse}
\end{equation}
where $\mathcal S$ is the simulator, $\mathcal H_j$ the measurement operator,
$y_j$ data, $s_j$ uncertainty/repeatability, and $\mathcal R$ a justified prior.
\textbf{Equation~\eqref{eq:inverse} is a proposed formulation, not an implemented
end-to-end inference engine.} Current functions fit the separate empirical
curves; unidentifiable parameter combinations must not be hidden by a low
training residual.

\subsection{Intended research-to-household response interface}
The implemented standalone response-table class provides bounded multilinear
interpolation on a full Cartesian grid. For a cell with lower/upper axis values
$x_j^-,x_j^+$, set $\lambda_j=(x_j-x_j^-)/(x_j^+-x_j^-)$:
\begin{equation}
 \widehat R(\mathbf{x})=\sum_{\mathbf b\in\{0,1\}^d}
 R_{\mathbf b}\prod_{j=1}^d
       \lambda_j^{b_j}(1-\lambda_j)^{1-b_j}. \label{eq:interpolation}
\end{equation}
Inside the cell the weights are nonnegative and sum to one. For example, a
hypothetical two-axis cell at $(\lambda_1,\lambda_2)=(1/4,1/2)$ has weights
$(3/8,1/8,3/8,1/8)$ under a compatible corner ordering. Interpolation is exact
for multilinear data, not for arbitrary nonlinear coupled material behavior.

The class validates grid completeness, units/physical bounds and evidence
provenance. It refuses absent/synthetic evidence and extrapolation by default;
explicit exploratory extrapolation is bounded to 10\% of each axis span and
still subject to response bounds. This mechanism does \emph{not} produce
validated evidence. No physical response tables are shipped, and neither native
mode currently consumes these tables. The dashed bridge in
Figure~\ref{fig:architecture} is intentionally not an implemented feedback loop.

A defensible future acceptance chain is: constitutive/balance tests; independent
space, timestep and domain-size studies; clean-surrogate measurements with
held-out errors no greater than 20\% or measured repeatability; then evidence-
qualified response reduction and measured household tests. Material variation,
parameter uncertainty, numerical error and behavior variability must be reported
separately. A confidence interval on a fitted parameter does not absorb unknown
model discrepancy.

\paragraph{Implementation mapping.}
\path{python/litter_physics/calibration.py}: \nolinkurl{fit_uptake},
\nolinkurl{fit_breakup}, \nolinkurl{fit_rheology}, \nolinkurl{friction_prior},
\nolinkurl{score_personalization}, \nolinkurl{apply_bundle}.
\path{python/litter_physics/response_tables.py}: \nolinkurl{ResponseTable},
\nolinkurl{Axis}, \nolinkurl{Evidence}, \nolinkurl{Evaluation}.
